\chapter{Introduction}

\section{Research background}

This project uses behaviour trees to control enemies in a game. A behaviour tree divides decisions into composite nodes, condition nodes and action nodes, so high-level strategy can be edited separately from movement, animation and combat code. For a small prototype, writing all logic in one script can still work. However, after states and interruption conditions increase, the control flow becomes increasingly difficult to check. The hierarchical structure of behaviour trees can reduce this problem, and this is also one of their main uses in game development (Colledanchise and Ögren, 2018; Iovino et al., 2022).

A node-link tree can directly show parent-child relationships, but the wider and deeper the tree becomes, the more canvas it occupies. Behaviour tree cards also show parameters, descriptions, Decorators, blackboard keys and running states, so they usually occupy more space than an ordinary hierarchy that only shows names. When many cards appear on the canvas, developers need to zoom, pan and search frequently, and the current execution path can also be drowned out by unrelated branches. This problem directly affects node location, condition modification and failure checking, not only whether the picture is tidy.

Therefore, this project first completed a Godot plugin that can save, run and debug real NPC decisions, and then used it to test the five display optimisation functions kept in the current version. These display optimisation functions respectively deal with dragging occlusion, information density during zooming, cards blocking connections, the relationship range of selected nodes, and local detail in an overview.

\section{Problem definition and research gap}

Existing research deals with large graphs from several directions. Tidy tree layout focuses on hierarchy, order and compactness, while layout adjustment research is more concerned with whether the original position and topology can be retained after overlap is removed (Reingold and Tilford, 1981; Sugiyama et al., 1981; Misue et al., 1995; Dwyer et al., 2006). Semantic zoom reduces displayed content when zooming out, while fisheye views enlarge the focus and keep surrounding structure (Bederson et al., 1996; Furnas, 1986). Interactive emphasis and EdgeLens further deal with related subgraphs and connection occlusion (Ware and Bobrow, 2005; Wong et al., 2003). These works show that large graphs can be improved from layout, information density, focus and connection treatment.

However, these studies do not directly explain how a game behaviour tree editor should combine these methods. A behaviour tree is different from an ordinary directory. Sibling order may decide execution priority, leaf nodes call Actor code, and Decorators may also prevent a branch from running. Therefore, a display optimisation function cannot only make the picture look tidier. It must also preserve execution order, support undo, and avoid breaking blackboard and runtime debugging.

\section{Research aim and question}

The aim of this dissertation is to design and evaluate composable large-tree display optimisation on a functionally verified Godot 4.6 visual behaviour tree platform that can drive real NPCs. The research question is: under different behaviour tree node scales and physical screen sizes, can the composable display optimisation proposed by this project improve the measurable display and editing interaction performance of large behaviour trees compared with a traditional baseline display, while keeping the behaviour tree structure and execution semantics unchanged?

The main contributions of this dissertation include three aspects. First, this project implemented a complete Godot 4.6 visual behaviour tree platform, so real game behaviour can be edited, run and debug in the same system. Second, the project designed and implemented five display optimisation functions that can be compared independently for the display problems of large behaviour trees. Finally, the project used five real game behaviour trees and three physical screen sizes for experiments, compared the effects and costs of different display optimisation functions in different environments, and introduced the project developer's structured usage evaluation, using this to explain which display optimisation functions are suitable as general functions and which are only suitable for specific use scenarios.

\section{Project objectives and success criteria}

The project work is divided into platform construction and research evaluation. Platform construction includes implementing serialisable tree, node, Decorator and typed blackboard resources; implementing the success, failure and running semantics of Root, Sequence, Selector, Reactive Selector, Random Selector, Parallel, Repeat, Wait, Action and Condition; allowing reusable components to assign a tree and Actor to an NPC, and allowing leaf nodes to call Actor methods; and providing creation, connection, disconnection, dragging, box selection, multi-node movement, typed editing, validation, saving, loading, undo and redo.

Before carrying out the display optimisation experiment, it is first necessary to confirm that the basic functions of the behaviour tree plugin are correct. Each type of node can execute as expected, the edited behaviour tree can be saved and reopened correctly, and undo and redo do not damage the tree content. Five playable enemies should respectively load the 31-, 61-, 121-, 241- and 364-node trees. The player has unlimited health, the number of enemies is limited, and the game can end after all enemies are defeated.

The success criteria of the display research are divided into three categories. First, a function must produce the direct change it claims, such as Smart Drag reducing overlap caused by dragging, Adaptive reducing overview card area, and Related Focus dimming unrelated nodes. Second, a function must not change tree topology, Decorator ownership or sibling execution order. Third, after a function is disabled, it must not leave size, opacity, border or temporary layout residue. If the test process has a program error or stops unexpectedly, the test is still counted as failed even if other checks pass.
